% Transcription par lecture visuelle directe, sans OCR.
% Statut : transcrit ; relecture indépendante à effectuer.
% Source : PDF original, vue 134, côté droit, page manuscrite 269, après le chapitre onzième.
\chapter{Troisième suite de la préfecture de M. Franchet. Chute de Napoléon.}

\section*{M. Franchet préfet}
% Mention marginale de la page manuscrite 269.
Napoléon après avoir échoué par l'issue du concile national, dans le projet qu'il avoit formé de séparer l'Église de son chef, et ne pouvant pas non plus réussir à vaincre à Savonne, ce qu'il appeloit l'obstination du Pape, se décida à le faire transférer dans le château de Fontainebleau. Il espéroit en le rapprochant ainsi de sa personne, parvenir enfin à triompher de sa résistance.

Le Saint père se plaignoit de n'avoir personne avec qui il pût conférer sur les concessions qu'on vouloit exiger de lui, puisqu'on lui avoit enlevé tous les Cardinaux. L'Empereur se détermina alors à lui rendre le Cardinal Pacca. Ce prince de l'Église avoit accompagné le Pape dans le voyage de Savonne ; mais à l'arrivée il en avoit été séparé et enfermé dans le château de Fenestrelle.
% Source : PDF original, vue 135, côté gauche, page manuscrite 270.
On le mit en liberté et on lui donna la permission d'aller à Fontainebleau joindre le Souverain Pontife ; c'étoit le vœu le plus ardent du Cardinal ; il se mit donc en route. Lyon étoit sur son passage ; il y arriva un vendredi soir et fut logé à l'hôtel de l'Europe ; il y rencontra M. le Vicomte de Montmorency.

\section*{Passage du Cardinal Pacca à Lyon}
% Mention marginale de la page manuscrite 270.
Les vertus du premier baron chrétien jetoient un si vif éclat, elles lui assuroient un tel ascendant dans la société, que le despote avoit craint l'influence qu'il pouvoit exercer dans la capitale ; il l'avoit exilé à quarante lieues de Paris. Nous fûmes assez heureux pour que M. de Montmorency voulût bien choisir notre ville pour le lieu de sa résidence ; c'est cette circonstance qui fut cause qu'il se trouva à l'hôtel de l'Europe où il s'étoit établi, lorsque le Cardinal Pacca y arriva.

Vous devez juger avec quel bonheur se virent ces deux illustres personnages, mais M. de Montmorency, fidèle associé de notre Congrégation, voulut en bon frère, partager avec nous la joie qu'il éprouvoit. Il fit prévenir le plus grand nombre de Congréganistes que la brièveté du délai put le lui permettre de la présence du Cardinal dans nos murs, et de son intention de dire le lendemain, avant de se remettre en route, la messe dans l'église de St François. Tous ceux qui purent être avertis se gardèrent bien de manquer à ce pieux rendez-vous, aussi la Congrégation s'y trouvat-elle presque au complet.

M. Guichard curé de la paroisse vint à la tête de son clergé, recevoir processionnellement le Cardinal à la porte de son église. Il lui adressa avec énergie un discours rempli de chaleur, par lequel il le félicitoit d'avoir souffert pour la cause de la religion et se félicitoit lui-même ainsi que son troupeau de pouvoir en ce
% Source : PDF original, vue 135, côté droit, page manuscrite 271.
jour révérer les marques de ces chaînes, qui avoient placé celui qui les avoit portées au nombre des confesseurs.

Pour s'exprimer ainsi sous le régime despotique auquel nous étions soumis, il falloit certainement avoir du courage ; heureusement la police de Napoléon n'avoit pas été avertie et ne se trouvoit pas là ; il n'en mésarriva pas au vénérable curé, et ce noble langage ne produisit d'autre effet que de remplir de joie les fidèles présents dont il exprimoit si bien les sentiments, et de verser dans l'âme du prince de l'Église auquel il étoit adressé, des consolations abondantes, car il dut alors juger, qu'il falloit que la foi fût encore bien vive dans un pays où sous le joug de fer d'un despote, au risque de tout ce qui pourroit en advenir, on ne craignoit pas de rendre ainsi un hommage public à la vérité.

Les membres de la Congrégation eurent ainsi qu'un grand nombre de personnes pieuses le bonheur de recevoir la communion des mains de son Éminence ; ce fut pour nous comme une espèce de communion générale. Comme les prières de cette pieuse assemblée pour la paix de l'Église étoient ferventes ! Comme la ferveur de tous se sentoit renouvelée en présence de ce digne représentant de l'Église militante ! Dieu lui-même avoit placé cette noble victime qu'il avoit choisie parmi les princes de son peuple sous les yeux de ses enfants chéris pour leur apprendre comment l'héroïsme de la résignation sait remporter la victoire malgré toute l'astuce et la perfidie des persécuteurs.

M. le V\up{te} de Montmorency avoit eu soin de faire souffler un petit mot à l'oreille de chaque Congréganiste, après la messe ils se trouvèrent réunis à l'hôtel de l'Europe, dans l'appartement qu'occupoit le Cardinal. Notre illustre associé eut
% Source : PDF original, vue 136, côté gauche, page manuscrite 272.
la bonté de les présenter tous à son Éminence ; il ne les désigna pas sous le titre de Congréganistes, il ne lui laissa pas ignorer que tous les Lyonnais qui se trouvoient présents étoient des hommes dévoués de cœur et d'âme à l'Église Catholique, Apostolique et Romaine. Toute cette pieuse réunion se prosterna aux pieds du Cardinal confesseur de la foi, et reçut avec émotion les bénédictions qu'il répandit sur elle, en éprouvant lui-même un vif attendrissement.

Il s'établit ensuite une conversation générale, dans laquelle le Prélat déploya une douce, affectueuse et spirituelle amabilité. Tout en parlant des circonstances où se trouvoit l'Église, on le félicitoit sur sa délivrance ; mais il étoit bien loin de la regarder comme complète. Voyez, nous disoit-il en souriant, tout en nous montrant un aigle qui placé au dessus de son lit en soutenoit les rideaux, voyez, il me tient sous sa griffe… Il ne se trompoit pas, il ne faisoit que changer de prison, car dès qu'il fut arrivé à Fontainebleau, il partagea celle du Pape.

L'Église gémissoit alors d'autant plus sur la captivité de son chef, qu'il devenoit de plus en plus difficile de communiquer avec lui. La police ne laissoit rien transpirer en dehors de ce qui pouvoit concerner le Père commun des fidèles. Les Catholiques étoient donc dans l'ignorance absolue de tout ce qui se passoit à Fontainebleau. Ils étoient réduits à prier pour le St Père, et à attendre le moment fortuné que la miséricorde Divine avoit marqué pour la délivrance.

Pendant que tous ceux qui avoient conservé la foi vive, éprouvoient si douloureusement cette absence totale de nouvelles d'un personnage qui leur étoit si cher, la Congrégation fut
% Source : PDF original, vue 136, côté droit, page manuscrite 273.
assez heureuse pour en recevoir de bien positives, grâce au Sauveur qui avoit bien voulu lui conserver un de ses associés. Vous devez vous rappeler que parmi les membres de la Congrégation de Paris que nous avions associés à celle de Lyon, se trouvoit cité M. de Rohan chabot. Depuis il étoit devenu chambellan de l'Empereur. Napoléon qui cherchoit à rattacher autant que possible l'ancienne noblesse à sa dynastie, saisissoit toutes les occasions d'appeler à sa cour ceux qui étoient porteurs d'un grand nom. C'est d'après ces principes qu'il avoit nommé M. de Rohan chabot chambellan de l'Empereur. Et lorsqu'il lui plaisoit ainsi de disposer de vous, il falloit nécessairement accepter, ou se résigner à l'exil. M. de Rohan avoit pris le premier parti.

Il n'avoit pas paru à Lyon depuis plusieurs années, lorsqu'un matin le Vice Préfet reçoit de lui un billet ainsi conçu : Mon cher ami, je n'ai que deux heures à passer à Lyon, j'ai les choses les plus intéressantes à vous communiquer ; venez, je vous prie de suite à l'hôtel de l'Europe où suis logé. Aussitôt le Vice Préfet y accourut, voici ce qu'il lui raconta. « Je partois pour l'Italie ; j'avois un grand désir de pouvoir recevoir en partant la bénédiction du Pape, et en même temps me mettre à ses ordres, pour tout ce dont il plairoit à Sa Sainteté de me charger pendant mon voyage. Il étoit inutile de demander une permission pour entrer au château, elle m'eût été très certainement refusée ; mais je savois que le gouverneur étoit absent ; je mets donc mon grand uniforme de chambellan dans ma valise, je monte dans une chaise de poste et je pars. Arrivé à Fontainebleau, je me revêts de mon habit d'ordonnance, je me présente à la grille du château et je demande à être introduit auprès du Saint Père. À la vue de mon uniforme, personne ne met en doute que je ne sois chargé de quelque message par l'Empereur,
% Source : PDF original, vue 137, côté gauche, page manuscrite 274.
aussi toutes les portes me sont elles ouvertes sans la moindre difficulté. J'ai eu la consolation de baiser les pieds du Saint père et de le supplier de me donner ses ordres et de répandre sur moi et sur les miens ses bénédictions. Le Pape m'a accueilli avec sa douceur et sa bonté accoutumée et m'a accordé bien des grâces particulières pour ma famille et pour moi. » M. de Rohan fit voir au vice Préfet un portrait de la Ste Vierge dont Sa Sainteté lui avoit fait présent. Je n'ai pas voulu, ajouta-t-il, passer à Lyon sans vous donner ces détails. Gardez les pour vous seul pendant deux jours, afin de me laisser le temps de m'éloigner avant que cela puisse s'ébruiter. Mais ensuite faites en part à tous nos amis que cela intéressera vivement. Vous pourrez leur dire en même temps que le Pape m'a paru en très bonne santé.

\section*{Nouvelles du Pape par M. de Rohan}
% Mention marginale de la page manuscrite 274 ; placée après le récit de sa visite.
Vous comprenez, Messieurs, avec quel bonheur le vice Préfet porta à chacune de nos sections, ces communications si précieuses pour les amis de la religion, et avec quelle joie elles furent accueillies par tous les Congréganistes. Outre l'intérêt qu'avoient par elles-mêmes pour nous les nouvelles du Saint père, cet incident offroit aux membres de la Congrégation la preuve que le bonheur d'appartenir à cette pieuse Société les plaçoit dans un rapport plus immédiat avec le père commun des fidèles que ses autres enfants. Cette pensée devenoit pour eux dans ces temps de persécution une source de consolations et un puissant motif d'encouragement.

La puissance de Napoléon étoit parvenue à son plus haut degré d'élévation ; il avoit subjugué l'Autriche ; il avoit subjugé la Prusse ; sous le titre de Protecteur de la Confédération du Rhin, il maîtrisoit le reste de l'Allemagne. Toute l'Italie et la Suisse étoient soumises à ses lois. Son frère Louis régnoit à Amsterdam,
% Source : PDF original, vue 137, côté droit, page manuscrite 275.
son frère Jérôme en Westphalie ; son beau frère Murat à Naples ; quoi que sur un trône bien chancelant son frère Joseph régnoit aussi à Madrid. La Pologne, docile à ses ordres, attendoit comme un oracle le son qui sortiroit de sa bouche impérieuse pour obéir à cette volonté souveraine en proclamant son roi. Souverain de la moitié de l'Europe, il arme les nations qu'il a vaincues et veut s'élancer avec elles sur les nations qui existent encore pour les subjuguer à leur tour. La plus violente commotion menace donc en ce moment d'ébranler le monde ; des armées innombrables se rassemblent de toutes parts pour aller sous la conduite du vainqueur de l'Europe, écraser l'Autocrate de toutes les Russies. Mais au milieu de la gloire qui l'environne, Napoléon s'irrite d'avoir rencontré à l'exécution de ses desseins, un obstacle qu'il n'avoit pas prévu, car l'orgueil, ce premier né de l'enfer est ainsi fait la plus mince résistance le jette dans un violent désespoir. Tout l'univers tomberoit aux pieds de ce soldat parvenu au faîte de la grandeur que le cœur de ce nouvel Aman ne sera pas satisfait tant que ce pauvre vieillard, nouveau Mardochée, persistera à ne pas courber sa tête vénérable sous le joug qu'il prétend lui imposer, tant qu'il refusera de reconnoître ce titre de roi de Rome, dont le despote pour enchaîner à jamais la ville des Césars à son char de triomphe, a salué son fils à sa naissance ; et tant que sa puissance spirituelle n’aura pas abaissé le pouvoir des clefs devant celui de l'aigle dévastateur. Aussi il n'attend plus que d'avoir vaincu les Russes, pour se débarrasser enfin de ce prêtre téméraire qui ose braver sa puissance. Le décret de déchéance est tout préparé, mais c'est de Moscou et de St Pétersbourg qu'il doit être daté, on assure même que pour conserver le souvenir de cette double
% Source : PDF original, vue 138, côté gauche, page manuscrite 276.
victoire, une médaille doit être frappée qui portera cet exergue : Dieu au ciel, Napoléon sur la terre.

Il faut convenir que ces sinistres projets avoient tout ce qu'il falloit pour exciter parmi les amis de la religion une anxiété bien vive ; cependant les Congréganistes étoient remplis de confiance, car ils se croyoient en droit d'espérer d'une espérance qui ne pouvoit pas être confondue, et voici sur quoi ils se fondoient ; dans la situation où se trouvoient alors les affaires, si Napoléon eut triomphé, et qu'après être demeuré maître du monde, il eut consommé le schisme qu'il méditoit, que fût devenue alors l'Église catholique ? Où auroit-elle pu trouver un refuge ? Nous nous étions appuyés sur les promesses de J. C. avec une conviction que rien ne pouvoit ébranler. L'Évangile est là. Le Seigneur a parlé. Le ciel et la terre passeront mais ses paroles ne passeront point. Il s'y est donc engagé, il faut qu'il se lève et qu'il venge sa cause. Comment cela arrivera-t-il ? nous l'ignorons ; mais au moment où nous nous y attendrons le moins, nous verrons une petite pierre se détacher de la montagne et renverser ce colosse qui en ce moment étonne l'Univers. Cette pensée remplissoit tous les cœurs Congréganistes ; elle soutenoit leur courage en présence de tous les maux dont nous pouvions être menacés.

\section*{Campagne de Russie}
% Mention marginale de la page manuscrite 276.
Cependant le grand homme toujours plein de confiance dans son étoile va prendre le commandement de ses armées et se croit sûr de la victoire. Il ne comprend pas que cette étoile prétendue qui semble jusqu'à ce jour l'avoir protégé n'est autre chose que la main du Seigneur Tout puissant, qui l'a conduit de succès en succès, jusqu'à ce que les desseins qu'il avoit sur lui, fussent accomplis. C'est cette main qui l'a
% Source : PDF original, vue 138, côté droit, page manuscrite 277.
dirigé pour rétablir la religion en France ; c'est encore cette main qui l'a armé de cette épée vengeresse, avec laquelle, nouveau fléau de Dieu, il a été chargé de punir l'Europe, comme la Révolution avoit été chargée de punir la France ; mais aujourd'hui cette mission est terminée ; il n'y a plus pour Napoléon de divin protecteur ; pour me servir de son langage son étoile va pâlir ; le despote a osé porter sur l'Oint du Seigneur une main sacrilège : son arrêt est prononcé.

L'Empereur traverse le Niémen à la tête de six cent mille hommes, pour ouvrir cette désastreuse campagne de Russie. La fatalité les entraîne, s'écrie-t-il avec arrogance ; que leurs destins s'accomplissent… Conquérant orgueilleux ! c'est ta propre bouche qui vient de prononcer ta condamnation. Oui, les destins, c'est à dire, les arrêts de la sagesse éternelle vont s'accomplir. Ne vois-tu pas la même main qui écrivit autrefois sur la muraille, la sentence de l'impie Balthasar ? Elle trace en ce moment la tienne. Marche ! marche ! L'Ange exterminateur est prêt. Voilà qu'un tourbillon de vent va emporter tes nombreuses armées comme une paille légère.

Toutefois la victoire n'abandonne pas encore ses aigles. Les Français s'avancent à pas de géant ; mais à mesure qu'ils pénètrent dans le cœur de l'empire Moscovite, un vaste désert se forme et grandit entre cette armée d'invasion qu'aucun obstacle ne peut arrêter et les magasins nécessaires à sa subsistance. Elle poursuit comme un torrent dévastateur sa course rapide jusqu'à Moscou. C'est là que la colère de [lacune]
% Source : PDF original, vue 139, côté gauche, page manuscrite 278.
l'attend ; il a inspiré aux Russes une grande pensée ; comme autrefois on vit des généraux se dévouer généreusement à une mort certaine pour sauver leurs concitoyens, aujourd'hui c'est une ville toute entière qui se sacrifie pour sauver la patrie. Moscou eut pu offrir à nos malheureux soldats un asile protecteur, Moscou n'est plus ! Un monceau de cendres, voilà tout ce qui nous reste au milieu du désert… Ce n'est pas encore tout ; à cette voix toute puissante qui appela autrefois le feu du ciel sur Sodôme et les eaux de la mer Rouge sur les Égyptiens, un vent glacial accourt du pôle du Nord. Il vient entourer nos malheureux compatriotes du froid rigoureux d'un hiver prématuré ; la plupart d'entre eux, tristes victimes de l'ambition de leur chef périssent de froid ou de faim, et servent de pâture aux oiseaux de proie ou aux ours du désert ; les autres devenus prisonniers des Russes, sont condamnés à aller gémir loin de leur patrie, dans cette province glacée qu'on appelle la Sibérie.

Napoléon signa alors ce bulletin fameux dans lequel les détails affreux des souffrances de ses compagnons d'arme, sont suivis de cette phrase d'un égoïsme barbare : Jamais l'Empereur n'a joui d'une meilleure santé. Il quitta ensuite son armée et retourne comme un fugitif dans la capitale de la France. Oh ! qui vous dira ce qui se passoit alors dans cette âme dont l'orgueil venoit d'être confondu d'une manière si éclatante ! ah ! sans doute, s'il ne l'avoue pas ouvertement, sa conscience doit en ce moment crier bien fort au fond de son cœur : Si Dieu est au ciel, il règne aussi sur la terre.

Bien loin d'être corrigé par le malheur et de profiter de la puissance qui lui reste pour donner enfin au monde la paix et le
% Source : PDF original, vue 139, côté droit, page manuscrite 279.
repos, Napoléon dans l'espoir de relever sa grandeur chancelante, impose à la France de nouveaux sacrifices ; il compte encore ressaisir la victoire. Vaine espérance ! elle lui a échappé pour toujours. Les peuples qu'il avoit subjugués saisissent tour à tour avec empressement l'occasion de secouer un joug si long et impatiemment supporté. Le voilà bientôt de toutes parts forcé de se réfugier dans les anciennes limites de la France, qu'il ne peut plus même préserver de l'invasion. Déjà le nord est envahi par les Russes et les Prussiens ; l'est par les Autrichiens et les troupes de la Confédération du Rhin ; le midi par les Anglais et les Espagnols. Napoléon ne cherche plus qu'à couvrir sa capitale ; c'est à la préserver de tomber entre les mains de la coalition que tendent maintenant tous ses efforts.

\section*{Délivrance du Pape}
% Mention marginale de la page manuscrite 279.
Pendant que le sceptre échappe déjà de cette main qui l'avoit porté avec tant de rigueur, la pensée du prisonnier de Fontainebleau se représente à l'imagination de l'Empereur. Peut-être alors éprouva-t-il comme Antiochus, cette espèce de contrition qui nous fait déplorer les offenses dont nous nous sommes rendus coupables envers la justice divine, à cause des malheurs qu'elles ont attirés sur nous. Peut-être aussi le même sentiment qui détermina jadis les Philistins à renvoyer l'arche Sainte dont la présence avoit appelé sur eux la colère du Seigneur, produisit-il en lui les mêmes effets. Quoiqu'il en soit, il se décida à rendre la liberté au chef de l'Église. Le voyage du Souverain Pontife fut un véritable triomphe ; il parcourut tout le midi de la France et une grande partie de l'Italie au milieu des témoignages de la vénération des peuples, et reprit sa place sur la chaire de St Pierre au bruit des acclamations unanimes de ses sujets ivres de joie d'avoir enfin retrouvé leur père.

C'est ainsi que deux fois en moins de quinze ans, la main du
% Source : PDF original, vue 140, côté gauche, page manuscrite 280.
Très haut est venue affermir l'Église ébranlée sur cette base qu'elle-même a posée. \textit{Tu es Petrus et super hanc petram, ædificabo Ecclesiam meam.}\footnote{(Note du transcripteur) Littéralement : « Tu es Pierre, et sur cette pierre je bâtirai mon Église ». Évangile selon Matthieu 16, 18.} Certes l'accomplissement des oracles de l'Évangile se montre ici avec une telle évidence, qu'il ne reste plus qu'à plaindre l'aveuglement des incrédules qui en présence de tant de merveilles ont le triste courage de persévérer dans leur incrédulité.

\section*{Abdication de Napoléon}
% Mention marginale de la page manuscrite 280.
Cependant Napoléon luttoit en vain contre la tempête dans laquelle sa puissance devoit s'anéantir. Son empire disparaissoit comme ces bulles de savon qu'un enfant fait éclore, qui se montrent un instant revêtues des plus brillantes couleurs et que le plus léger souffle fait évanouir. Déjà il se trouvoit réduit à quelques misérables lambeaux que chaque jour on arrachoit pièce à pièce. Pressé de tous côtés par les nations coalisées contre lui, il voit tomber jusqu'au dernier tous les fleurons de cette couronne impériale qui naguère faisoit sécher de frayeur et les peuples et les rois. Voilà sa capitale envahie. Il entend les cris de joie et d'amour des Français qui rendus à la liberté demandent de toute part à rentrer sous la domination de leurs souverains légitimes. Enfin, ô justice de mon Dieu ! c'est dans ce même château où il a retenu prisonnier le vicaire de J. C. dont il vouloit usurper l'autorité, qu'il est réduit à signer son abdication, et c'est de là qu'il part pour aller à l'île d'Elbe où il est confiné pour un temps, jusqu'à ce qu'il plaise à la Providence de le rappeler encore pour châtier de nouveau la France et le reléguer ensuite sur un rocher, dans la vaste étendue des mers. C'est là qu'il a pu comprendre en méditant sur sa destinée que malgré tous les rêves pompeux
% Source : PDF original, vue 140, côté droit, page manuscrite 281.
que l'orgueil de l'homme peut concevoir, il n'est rien et ne peut rien contre les suprêmes volontés de la souveraine majesté, qui étoit hier, qui est aujourd'hui, qui sera demain et toujours dans les siècles des siècles.

Nous avons dû, Messieurs, en écrivant cette histoire, raconter avec un certain développement la chute de Napoléon. Quoique ce grand événement appartienne à l'histoire de France, et que d'ailleurs vous en ayez tous une parfaite connoissance, il n'étoit pas moins nécessaire de lui donner une place dans la suite de notre récit, car cette chute étoit une délivrance pour la Congrégation, en même temps qu'elle en étoit une pour l'Église et pour la France. Cela nous a donné d'ailleurs une nouvelle occasion d'admirer avec quelle sagesse la divine providence dispose toujours les événements pour la conservation de son Église et le plus grand bien de ses élus.

Portons de nouveau nos regards sur la Congrégation ; voyons ce qu'elle étoit devenue au moment où cette grande Révolution alloit apporter des changements si notables dans la situation de notre patrie. Vous avez déjà remarqué dans le chapitre précédent qu'à mesure que la position de Napoléon devenoit plus critique, la Congrégation étoit de plus en plus oubliée. Aussi plusieurs mois avant que les alliés se rendissent maîtres de Lyon, ce qui eut lieu trois semaines avant l'occupation de Paris, personne au monde n'y pensoit plus et notre sécurité étoit complète.

Cependant, il faut le dire, les sections n'en étoient pas plus actives. Depuis que la guerre s'étoit rapprochée de nous, et bien plus encore depuis l'occupation de la ville, trop de soins exigeoient de la part des membres une attention de tous les instants, pour qu'on pût espérer d'eux une certaine exactitude à se rendre aux assemblées. Une seule assemblée générale avoit eu lieu (celle dont nous avons
% Source : PDF original, vue 141, côté gauche, page manuscrite 282, avant le chapitre treizième.
déjà parlé) pour l'élection du Préfet. Sans être dans un état complet de mort nous agissions donc fort peu. Les Congréganistes n'en étoient pas moins toujours sincèrement dévoués à la Congrégation et à tout le bien qu'elle s'y proposoit d'opérer en se réunissant. Plus que jamais, ils se livroient à l'espoir de voir luire bientôt des jours meilleurs ; ils attendoient avec confiance le moment qui ne paroissoit plus éloigné où la Congrégation pourroit retrouver toute la vigueur de ses premiers ans.

Cette heureuse époque fut celle où le Roi légitime fut enfin replacé sur le trône de ses ancêtres. Cet événement fut un grand sujet de joie pour les Congréganistes, comme chrétiens, comme Français et aussi dans l'intérêt de notre société ; car ils y voyoient en même temps la liberté de nos frères captifs assurée et la facilité de rétablir dans tout leur ensemble les rapports si doux qui nous unissoient. On peut dire que la Congrégation ressembloit alors à une montre qui n'est pas montée ; pour lui rendre le mouvement et la vie, il suffit de remplir cette petite formalité. Pour les rendre à la Congrégation, il suffira de remettre son règlement en pleine activité.

Nous allons dans le chapitre suivant mettre sous vos yeux cette heureuse résurrection.
